\section{Методы компиляции в WebAssembly}

    \quad WebAssembly как целевая платформа компиляции поддерживается широким набором инструментариев, различающихся по исходным языкам, архитектуре компилятора и характеристикам генерируемого кода. Выбор конкретного стека определяет архитектуру бинарного модуля, объём сопутствующего рантайма и методику взаимодействия с хост-средой. Рассмотрим основные методы компиляции.

    \quad \textbf{AOT-компиляция (ahead-of-time).} Данный метод предполагает предварительную трансляцию исходного кода непосредственно в бинарный модуль WebAssembly ещё до запуска программы. Именно AOT-компиляция стала основным способом доставки производительных Wasm-модулей в браузер и серверные рантаймы, поскольку она обеспечивает предсказуемое время инициализации, компактный формат поставки и высокий уровень оптимизации. К инструментам этого класса относятся Emscripten, LLVM/Clang, WASI SDK, Rust с \texttt{wasm-bindgen} и \texttt{wasm-pack}, TinyGo, AssemblyScript и другие компиляторные цепочки, генерирующие Wasm напрямую.

\begin{table}[htbp]
\centering
\scriptsize
\renewcommand{\arraystretch}{1.15}
\begin{tabular}{|p{5.8cm}|p{5.8cm}|}
\hline
\textbf{Плюсы} & \textbf{Минусы} \\
\hline
Высокая производительность исполнения & Требуется отдельный этап сборки \\
\hline
Глубокие оптимизации на этапе компиляции & Плохо подходит для динамической генерации кода \\
\hline
Хороший контроль над экспортируемым интерфейсом & Может требовать сложной настройки памяти и рантайма \\
\hline
Удобная интеграция в стандартные сборочные пайплайны & Работа с внешними зависимостями бывает нетривиальной \\
\hline
\end{tabular}
\caption{Преимущества и недостатки AOT-компиляции}
\label{tab:aot_pros_cons}
\end{table}

    \quad \textbf{JIT-компиляция (just-in-time).} При данном методе код компилируется во время исполнения, когда уже известен контекст запуска, профиль нагрузки и характеристики среды. В классическом виде JIT чаще применяется не как основной способ получения Wasm-модуля из исходного текста, а как механизм динамической компиляции промежуточного представления или байткода внутри среды выполнения. В экосистеме WebAssembly этот путь встречается заметно реже, чем AOT, однако важен для рантаймов, где требуется адаптация к данным профилирования и возможность поздней оптимизации. Характерные инструменты и движки этого класса --- Wasmer JIT, Cranelift в JIT-сценариях, а также внутренние JIT-пайплайны движков V8 и SpiderMonkey.

\begin{table}[htbp]
\centering
\scriptsize
\renewcommand{\arraystretch}{1.15}
\begin{tabular}{|p{5.8cm}|p{5.8cm}|}
\hline
\textbf{Плюсы} & \textbf{Минусы} \\
\hline
Оптимизация с учётом реального профиля выполнения & Высокая сложность реализации \\
\hline
Возможность улучшать горячие участки после запуска & Увеличивает время холодного старта \\
\hline
Гибкая адаптация к среде исполнения & Требует более сложного рантайма \\
\hline
Подходит для сценариев поздней оптимизации & Редко используется как основной метод доставки Wasm \\
\hline
\end{tabular}
\caption{Преимущества и недостатки JIT-компиляции}
\label{tab:jit_pros_cons}
\end{table}

    \quad \textbf{Binary translation.} Этот метод основан на переводе уже существующего бинарного кода, байткода или инструкций другой архитектуры в представление, пригодное для исполнения в среде WebAssembly. Такой подход применяется тогда, когда исходный код недоступен либо перенос системы с уровня исходников экономически нецелесообразен. В практических и исследовательских решениях сюда можно отнести CheerpX, некоторые QEMU/LLVM-ориентированные схемы динамической трансляции и другие специализированные системы двоичного переноса.

\begin{table}[htbp]
\centering
\scriptsize
\renewcommand{\arraystretch}{1.15}
\begin{tabular}{|p{5.8cm}|p{5.8cm}|}
\hline
\textbf{Плюсы} & \textbf{Минусы} \\
\hline
Позволяет использовать уже существующий бинарный код & Очень высокая сложность реализации \\
\hline
Не всегда требует доступа к исходникам & Производительность обычно ниже AOT-компиляции \\
\hline
Полезен для переноса legacy-систем & Результат зависит от качества сопоставления архитектур \\
\hline
Может сокращать затраты на переписывание системы & Используется в основном в нишевых сценариях \\
\hline
\end{tabular}
\caption{Преимущества и недостатки binary translation}
\label{tab:binary_translation_pros_cons}
\end{table}

    \quad \textbf{Интерпретация.} В интерпретационном подходе в WebAssembly переносится не сама программа в виде нативно скомпилированного модуля, а интерпретатор, который затем исполняет внешний бинарный код, байткод или сценарный язык инструкция за инструкцией. Этот метод особенно удобен для обеспечения совместимости и быстрого запуска существующих экосистем без радикальной переработки их внутренней модели исполнения. К подобным решениям относятся Pyodide с интерпретатором CPython, Lua- и QuickJS-рантаймы, а также различные эмуляторы и интерпретаторы, скомпилированные в Wasm.

\begin{table}[htbp]
\centering
\scriptsize
\renewcommand{\arraystretch}{1.15}
\begin{tabular}{|p{5.8cm}|p{5.8cm}|}
\hline
\textbf{Плюсы} & \textbf{Минусы} \\
\hline
Высокая гибкость & Производительность ниже AOT-решений \\
\hline
Сравнительно быстрый перенос сложных экосистем & Часто уступает и JIT-подходам \\
\hline
Удобна для задач совместимости & Каждая инструкция проходит дополнительный уровень обработки \\
\hline
Позволяет сохранить полноту поддержки языка & Плохо подходит для тяжёлых вычислительных задач \\
\hline
\end{tabular}
\caption{Преимущества и недостатки интерпретации}
\label{tab:interpretation_pros_cons}
\end{table}

    \quad \textbf{Hybrid-подход.} Гибридная схема объединяет несколько методов компиляции и исполнения в рамках одной системы. На практике это означает, что часть кода может исполняться через интерпретацию, часть --- быть заранее скомпилированной AOT, а наиболее критические участки --- дополнительно оптимизироваться специальным рантаймом. Подобные решения характерны прежде всего для управляемых платформ и больших фреймворков, где необходимо одновременно сохранить совместимость, приемлемую скорость запуска и достаточную производительность. Типичными примерами являются Blazor WebAssembly с AOT-режимом, смешанные пайплайны Mono и ряд managed runtime-решений для Java и .NET.

\begin{table}[htbp]
\centering
\scriptsize
\renewcommand{\arraystretch}{1.15}
\begin{tabular}{|p{5.8cm}|p{5.8cm}|}
\hline
\textbf{Плюсы} & \textbf{Минусы} \\
\hline
Гибкий баланс между скоростью запуска и производительностью & Архитектурно более сложен \\
\hline
Позволяет сочетать сильные стороны разных схем исполнения & Требует большего объёма рантайма \\
\hline
Удобен для крупных управляемых платформ & Сложнее в отладке и сопровождении \\
\hline
Помогает сохранять совместимость платформы & Итоговое поведение зависит сразу от нескольких механизмов \\
\hline
\end{tabular}
\caption{Преимущества и недостатки hybrid-подхода}
\label{tab:hybrid_pros_cons}
\end{table}

    \quad \textbf{VM $\rightarrow$ Wasm.} При данном методе в WebAssembly компилируется виртуальная машина или рантайм, а уже внутри него исполняется байткод исходной платформы. В отличие от более общей интерпретации, здесь обычно речь идёт о полноценных управляемых экосистемах с собственным форматом байткода, системой типов, механизмами загрузки классов или сборки мусора. Этот путь особенно распространён для Java- и .NET-платформ, где требуется сохранить существующий стек библиотек и модель исполнения. Наиболее показательные примеры --- Blazor WebAssembly с рантаймом Mono, Java-ориентированные решения вроде CheerpJ и другие системы, в которых виртуальная машина переносится в браузер как Wasm-модуль.

\begin{table}[htbp]
\centering
\scriptsize
\renewcommand{\arraystretch}{1.15}
\begin{tabular}{|p{5.8cm}|p{5.8cm}|}
\hline
\textbf{Плюсы} & \textbf{Минусы} \\
\hline
Высокая совместимость с исходной экосистемой & Увеличивает размер поставляемого модуля \\
\hline
Можно повторно использовать существующий байткод & Даёт более медленный холодный старт \\
\hline
Сохраняет привычную модель исполнения платформы & Создаёт дополнительные накладные расходы рантайма \\
\hline
Облегчает перенос существующих библиотек & Эффективность зависит от качества реализации виртуальной машины \\
\hline
\end{tabular}
\caption{Преимущества и недостатки метода VM $\rightarrow$ Wasm}
\label{tab:vm_to_wasm_pros_cons}
\end{table}

    \quad Сопоставление перечисленных методов показывает, что между ними существует устойчивый компромисс. Чем ближе метод к прямой предварительной компиляции, тем выше, как правило, производительность и тем меньше объём вспомогательного рантайма. Чем сильнее система опирается на интерпретацию, виртуальную машину или гибридную модель, тем проще сохранить совместимость с исходной платформой, но тем выше накладные расходы на запуск и исполнение. Сравнительные характеристики основных методов приведены в таблице~\ref{tab:compilation_methods}.

\begin{table}[htbp]
\centering
\scriptsize
\renewcommand{\arraystretch}{1.2}
\begin{tabular}{|p{2.4cm}|p{2.0cm}|p{2.8cm}|p{2.3cm}|p{3.5cm}|}
\hline
\textbf{Метод} & \textbf{Вход} & \textbf{Производительность} & \textbf{Сложность} & \textbf{Использование} \\
\hline
AOT & Исходный код & Высокая & Средняя & Основной способ сборки Wasm-модулей \\
\hline
JIT & Код во время выполнения, IR, байткод & Средняя & Высокая & Редко, в основном во внутренних рантаймах \\
\hline
Binary translation & Готовый бинарь, машинный код, байткод & Средняя или низкая & Очень высокая & Нишевые сценарии переноса legacy-систем \\
\hline
Интерпретация & Бинарь, байткод, сценарный код & Низкая & Низкая или средняя & Совместимость и быстрый перенос экосистем \\
\hline
Hybrid & Смешанный вход & Средняя & Высокая & Managed runtime и крупные платформы \\
\hline
VM $\rightarrow$ Wasm & Байткод виртуальной машины & Средняя & Средняя & Java-, .NET- и другие байткодные экосистемы \\
\hline
\end{tabular}
\caption{Сравнение методов компиляции в WebAssembly}
\label{tab:compilation_methods}
\end{table}
